\subsection{Video Personalization}
\label{sec:video_personalization}

Video personalization aims to generate videos that maintain consistent identity with the user-provided reference~\citep{dreamvideo, li2024personalvideo, consisid}.
In this section, we integrate advanced personalization techniques into the video generation pipeline, achieving state-of-the-art performance.
The subsequent discussion provides a comprehensive elaboration of our approach, encompassing the model architecture, personalization data, and experimental validation.


\subsubsection{Model Design}

The core techniques behind video personalization tackle two primary challenges: (1) the acquisition of high-fidelity personalized identities and (2) the seamless integration of these identity features into the video generation pipeline.

We start from our \method-T2V foundation model by first obtaining personalized identity information. 
Existing video personalization approaches usually rely on ID extractors (\textit{e.g.,} ArcFace~\citep{deng2019arcface}) or versatile visual extractors (\textit{e.g.,} CLIP~\citep{radford2021learning}) to obtain identity information.
While they make promising progress, their performance is bottlenecked by the limitations of the feature extractors. 
For instance, ID extractors mainly focus on features for face recognition but may not capture other critical visual cues such as scars or stickers. 
Moreover, they are susceptible to low-resolution faces, illumination variations, and partial occlusions (such as glasses, masks, or hair). 
On the other hand, versatile visual extractors are not specifically tailored to the face domain and tend to focus on coarse-grained semantic information. 
Therefore, we choose not to rely on any feature extractor to avoid information loss, and our generation process is directly conditioned on the input face image in the latent space of \method-VAE.

There are many ways to inject the personalized identity information into the video generation process. In our experiments, we found that cross-attention operation is suited to interact with dense representations obtained from identity extractors, while self-attention is more appropriate for modeling data in the same latent space. Fig.~\ref{fig:framework_personalization} illustrates the design of our video personalization approach, which follows a self-attention paradigm.
Specifically, in the latent space, we first extend $K$ additional frames prepended to the given video using the segmented face images. 
These face images are extracted from the paired video using face detection and segmentation, and the facial landmarks are further aligned to a black canvas that has the same size as the video frame. 
Then, along the channel dimension of this extended video, we concatenate the face images with all-ones masks for the first $K$ frames, and blank images with all-zeros masks for the remaining frames. 
These concatenated signals serve as the conditions. 
Finally, the diffusion process is performed on this temporally extended video, conditioned on the channel-wise condition signals in an inpainting fashion. 
Note that no other modifications are applied to our \method-T2V model. 

During training, we randomly drop a portion of face images in the $K$ extended frames, so as to support $0$ to $K$ reference face(s) video generation. 
In the inference stage, starting from the pure noise of the extended video, we concatenate the user-provided face images (no more than $K$) to the extended frames in the channel dimension and set the corresponding masks as all-ones. 
The goal is to reconstruct the face images in the first $K$ frames, and synthesize a new personalized video in the subsequent frames that preserves the provided identity.



\begin{figure}[t]
    \scriptsize
    \centering
    \includegraphics[width=1.0\linewidth]{figures/personalization/video_personalization.pdf}
    \caption{The core design of the video personalization approach in the latent space of \method-VAE.}
    \label{fig:framework_personalization}
\end{figure}



\subsubsection{Dataset}
We curate a collection of personalization data from our T2V foundation model's training dataset through careful filtering and automatic data synthesizing. 
We first filter out about $\mathcal{O}(100)$M videos using our internal human classifiers and perform face detection at 1 FPS on these videos. 
A video shall be discarded if more than one face is detected in any frame or if over 10\% of the frames have no face detected. 
Then we calculate ArcFace similarities between consecutive frames and discard videos with low similarities. 
Afterwards, we perform face segmentation to remove the distracting backgrounds, and perform face landmark detection so as to facilitate canvas alignment during our training. 
Note that we do not filter out faces with the small areas, because we find that such videos usually contain full-body human figures. 
Finally, we construct about $\mathcal{O}(10)$M personalized videos, where each video is accompanied by 5 segmented faces in average.

We also resort to automatic data synthesizing to improve facial diversity. 
Particularly, we randomly select $\mathcal{O}(1)$M personalized videos from above, and adopt Instant-ID~\citep{wang2024instantid} to synthesize diverse faces for each video. 
Notably, we build a text template with over 100 prompts, including anime, line art, cinematic, Minecraft, and so on. Every time we take a random prompt from this template, as well as a random human pose estimated from the rest of the videos, to serve as inputs to Instant-ID. We further measure the ArcFace similarity on the generated faces and filter out those with lower similarities. 
In total, we gather about $\mathcal{O}(1)$M videos with synthesized faces, which greatly improve the diversity of styles, poses, illuminations, and occlusions on our personalization dataset.


\begin{table}[!h]
    \centering
    \begin{tabular}{l| c c c c c}
        % \toprule
         & \method & CN-TopA & CN-TopB & CN-TopC \\ %\midrule
        \shline
        \rule{0pt}{15pt} % 手动设置行高
        % Visual Quality
        Arcface Similarity & 
            0.5526 & 
            0.5655 & 
            0.5197 &
            0.4998\\

    \end{tabular}
    \caption{Comparison of other leading video personalization methods.}
    \label{tab:video_personalization_compare}
\end{table}



\subsubsection{Evaluation}
We present our video personalization results in Fig.~\ref{fig:ID_cases}, where the left part represents the input identities and the right part represents the synthesized personalized videos. 
All the testing images are randomly picked from Pexels\footnote{www.pexels.com}. 
We also evaluate our model on an unseen evaluation set and measure the ArcFace similarity between input faces and output personalized videos by sampling faces at 1 FPS from videos. 
The final similarity scores are averaged from these sampled faces. 
The corresponding results are illustrated in Table~\ref{tab:video_personalization_compare}, where our approach demonstrates competitive video personalization performance compared to other leading commercial and closed-source Chinese competitors.

\begin{figure}[htbp]
    \scriptsize
    \centering
    \includegraphics[width=\linewidth]{figures/personalization/ID_cases.pdf}
    \caption{Visualization of our video personalization approach.}
    \label{fig:ID_cases}
\end{figure}


